% ═══════════════════════════════════════════════════════════════════════════
% Chapter 17: Related Work and Comparative Analysis
% ═══════════════════════════════════════════════════════════════════════════

\chapter{Related Work and Comparative Analysis} \label{chap:related_work}

\section{Overview}

This chapter surveys related academic and commercial work in IMU-based motion tracking, VBT measurement systems, and hardware acceleration of sensor fusion algorithms. The analysis identifies gaps this dissertation addresses.

\section{Academic Literature}

\subsection{IMU-Based Barbell Tracking}

\subsubsection{Perez-Castilla et al. (2019)}

Found Limits of Agreement +/-0.11 m/s against linear position transducer. Devices were commercial black boxes; algorithm transparency limited.

\textbf{Gap addressed:} Open-source pipeline with reproducible algorithms.

\subsubsection{Banyard et al. (2017)}

Reported correlation r = 0.94 with LPT criterion. However, systematic bias of -0.05 m/s observed at high velocities.

\textbf{Gap addressed:} Sub-0.07 m/s MAE with VQF-based correction.

\subsubsection{Laidig et al. (2021)}

Established VQF as state-of-the-art for 6-DOF tracking. This work extends VQF to VBT-specific requirements (ZUPT, per-rep smoothing).

\textbf{Contribution:} First application of VQF to weightlifting; validated on 622 reps.

\section{Commercial VBT Devices}

\subsection{Comparative Summary}

\begin{table}[h]
    \centering
    \caption{Complete system comparison}
    \label{tab:system_comparison}
    \begin{tabular}{lccccc}
        \toprule
        \textbf{System} & \textbf{Accuracy} & \textbf{Latency} & \textbf{Price} & \textbf{Open} \\
        \midrule
        PUSH Band 2.0 & 0.10 m/s & 100 ms & \$299 & No \\
        Vmaxpro & 0.08 m/s & 80 ms & \$399 & No \\
        GymAware (LPT) & 0.02 m/s & 20 ms & \$799 & No \\
        Beast Sensor & Unknown & Unknown & \$149 & No \\
        \midrule
        \textbf{This work} & \textbf{0.068 m/s} & \textbf{88 ms} & \textbf{\$150} & \textbf{Yes} \\
        \bottomrule
    \end{tabular}
\end{table}

\section{Thesis Differentiation}

This dissertation contributes:
\begin{enumerate}
    \item Accuracy comparable to best commercial devices (0.068 m/s MAE)
    \item Full open-source release of hardware, firmware, algorithms
    \item Comprehensive validation (622 annotated repetitions)
    \item Complete FPGA/ASIC migration path
    \item Projected \$150 BOM vs. \$300-400 commercial pricing
\end{enumerate}

\section{Summary}

Related work demonstrates growing interest in IMU-based VBT but leaves gaps in transparency, validation rigor, and hardware acceleration. This dissertation addresses these gaps.

\newpage